\chapter{Regionalism, Decentralisation and the Question of Minority}

\section{Democratic Decentralisation: Gram Swaraj and Panchayati Raj}

\begin{KeyConcept}
\begin{itemize}
\item \textbf{Gandhi's Gram Swaraj} --- he distrusted centralized politics ('Raj Niti', where people are mere voters) and proposed decentralising power to the village ('Lok Niti'/'Lok Shakti'). The \textbf{ideal village} would be self-sufficient (in food and cotton), with a cattle reserve, playgrounds, water works and sanitation; those with wealth would act as \textbf{trustees}. National elites would attack caste/gender barriers while grassroots leaders would demand their implementation.
\item \textbf{Ambedkar's critique}: the Indian village is '\textbf{a sink of localism, a den of ignorance, narrow-mindedness and communalism}' --- Gram Swaraj would \textbf{reproduce caste} (the dominant caste would rule and exploit Dalits and women).
\item The \textbf{normative model of decentralisation} needs both \textbf{institutions} (representation of interest groups, people's participation, political/financial autonomy, experts/planners, integration with state/centre) and a conducive \textbf{environment}. This was given constitutional status by the \textbf{73rd and 74th Amendments} --- three tiers (village, block, district), with reservation for women, SCs and STs.
\item The \textbf{three Fs --- Funds, Functions, Functionaries} --- determine success: \textbf{Kerala} devolved these (a third of funds directly to Panchayats) and performs best; Bihar did not. \textbf{Problems}: \textbf{surrogate leadership} ('Sarpanchpati' --- the husband ruling in the wife's name), the state not devolving power, no experts, symbolic integration.
\item \textbf{Positive studies}: \textbf{Anandi Mani et al.} --- where women head Panchayats, reporting of crimes against women rises and women's issues (drinking water) are prioritised; \textbf{Raghabendra Chattopadhyay \& Esther Duflo} --- female Pradhans invest more in drinking water/sanitation, SC Pradhans in housing for SC areas.
\end{itemize}
\end{KeyConcept}

\begin{QuickRevision}
\begin{itemize}
\item Gandhi's Gram Swaraj (ideal village, trustees, Ramaraj) vs Ambedkar ('sink of localism' --- reproduces caste).
\item 73rd/74th Amendments (3 tiers; reservation); three Fs (Funds/Functions/Functionaries); Kerala best; problems = surrogate leadership (Sarpanchpati).
\item Mani et al. (female Sarpanch $\rightarrow$ reporting up), Chattopadhyay \& Duflo (female $\rightarrow$ water/sanitation; SC $\rightarrow$ housing).
\end{itemize}
\end{QuickRevision}

\section{Regionalism}

\begin{KeyConcept}
\begin{itemize}
\item A \textbf{region} is defined administratively or culturally; \textbf{regionalism is a cultural phenomenon} --- it emerges from \textbf{over-pride in one's culture} or from \textbf{relative deprivation} (e.g. the demands for Uttarakhand and Jharkhand). \textbf{Rasiduddin Khan} gave ten criteria to identify a region: language/dialect, social composition, ethnic regions, demographic features, geographic contiguity, cultural patterns, economy, historical antecedents, political background, and psychological make-up.
\item \textbf{Three forms}: \textbf{supra-state} (the South vs North on Hindi imposition), \textbf{inter-state} (water/border disputes; location of central projects/IITs), and \textbf{intra-state} (Gorkhaland).
\item Regionalism creates \textbf{narrow identities} --- the \textbf{'son of the soil'} movement (Shiv Sena in Maharashtra, the anti-migrant movement in Assam) claiming the first right over local resources; the \textbf{Khalistan} movement was the only genuinely secessionist one.
\item \textbf{But regionalism is also positive}: it counters centralisation, preserves linguistic/cultural identities, raises regional issues nationally, and --- as \textbf{Rajni Kothari} says --- makes Indian democracy more vibrant (federalism is a \textbf{compromise between nationalism and regionalism}). When linguistic states were created, many feared the \textbf{'balkanisation'} of India --- but the linguistically-created states (the South) performed \textbf{better} on socio-economic indicators, so India should follow a \textbf{plural integration model}, not a unitary one.
\end{itemize}
\end{KeyConcept}

\begin{QuickRevision}
\begin{itemize}
\item Region: administrative vs cultural; regionalism from over-pride or relative deprivation (Uttarakhand/Jharkhand); Rasiduddin Khan's 10 criteria.
\item Three forms: supra (South vs North), inter (water/border), intra (Gorkhaland); son of the soil (Shiv Sena); Khalistan = secessionist.
\item Positive: counters centralisation, preserves culture, Kothari (vibrant democracy), plural integration; linguistic states performed better.
\end{itemize}
\end{QuickRevision}

\section{The Question of Minority}

\begin{KeyConcept}
\begin{itemize}
\item The concept of \textbf{minority} arises from \textbf{urbanisation} (secondary urbanisation --- groups that do not follow the dominant culture) and from the \textbf{nation-state} (communities that do not constitute a nation). Legally, India recognises \textbf{religious and linguistic} minorities.
\item \textbf{Muslims} are the central minority concern: the \textbf{Sachar Committee} found them \textbf{worse off than Dalits}. Reasons for their backwardness: the \textbf{abolition of zamindari} (many Muslims were dependants), the \textbf{destruction of handicraft/artisan industries}, not taking \textbf{modern education} (preferring madarsas), the \textbf{migration of the educated to Pakistan}, and \textbf{discrimination in employment/housing}. They also face \textbf{communal violence and dehumanisation} (stereotypes about meat-eating, 'love jihad', the Ahmedabad non-veg rule).
\item Minorities have internal problems too --- \textbf{caste} (among Sikhs) and \textbf{sectarian} divisions (Shia--Sunni, Ahmadiyya). On Muslim women, the criminalisation of triple talaq is welcomed as deterrence/empowerment, but it raises the demand for a \textbf{Uniform Civil Code} --- which should be based on modern principles of gender equality, not the Hindu code. \textbf{Hilal Ahmed} notes that India is unique in the visibility of its Muslim minority (mosques, dress) --- which is now being turned into the basis of an exclusionary nationalism.
\item \textbf{Naxalism} (sociologically) is a product of \textbf{differential development} --- states/communities denied the benefits of modern development (exploited earlier by caste/class, now by state policies) imagine a different India; they are \textbf{anti-state} (attacking state symbols, not ordinary people). It began in Punjab (Gulzar's film \emph{Machis}) but dissipated there after the Green Revolution. The \textbf{farmers' protest} also created a new \textbf{'township of protest'} (makeshift homes, facilities, more women and Dalit participation).
\end{itemize}
\end{KeyConcept}

\begin{QuickRevision}
\begin{itemize}
\item Minority from urbanisation + nation-state; religious/linguistic; Muslims (Sachar --- worse than Dalits; zamindari abolition, handicraft destruction, madarsas, migration, discrimination, communal violence, dehumanisation).
\item Uniform Civil Code (modern, not Hindu); Hilal Ahmed (visibility of Muslims); caste/sect within minorities.
\item Naxalism = differential development, anti-state; farmers' protest = 'township of protest'.
\end{itemize}
\end{QuickRevision}
